\chapter{A Concrete Naming Certificate for $\BishopCauchyModulusEnvelopeUp$}
\label{ch:concrete-instances-bishop-cauchy-modulus-envelope-namecert}
\origin{human}

Following Bishop--Bridges constructive analysis, the $\BishopCauchyModulusEnvelopeUp$ packet names the finite envelope through which a Cauchy-modulus request is routed before the L10 Real and Completion spine reads a terminal seal. It sits beside the finite observation discipline of \autoref{ch:concrete-instances-streamname-namecert}, the regular rational readback of \autoref{ch:concrete-instances-regseqrat-namecert}, the dyadic tolerance core of \autoref{ch:concrete-instances-dyadicratcore-namecert}, the directed modulus comparison surface of \autoref{ch:concrete-instances-dyadic-cauchy-modulus-lattice-namecert}, and the terminal real-name seal of \autoref{ch:concrete-instances-real-namecert}. The packet binds those sibling routes into one displayed finite envelope; it is not a selected limit, quotient stream equality principle, countable-choice operation, host modulus function, or ambient completeness claim for $\RealUp$.

\begin{definition}[Bishop Cauchy modulus envelope carrier]
\label{def:bishop-cauchy-modulus-envelope-carrier}
A $\BishopCauchyModulusEnvelopeUp$ carrier is a finite $\BHist$ packet
$$
E=(M,S,R,D,J,A,H,C,P,N)
$$
with the following displayed rows:
\begin{enumerate}
\item $M$ is the modulus request row, naming the requested dyadic precision and the finite tail window demanded by a Bishop-style Cauchy argument.
\item $S$ is the $\StreamNameUp$ finite-window row opened exactly at the indices exposed by $M$.
\item $R$ is the $\RegSeqRatUp$ readback row obtained from those stream windows.
\item $D$ is the $\DyadicRatCoreUp$ tolerance ledger recording the finite dyadic comparisons used by the request.
\item $J$ is the $\DyadicCauchyModulusLatticeUp$ comparison row, placing the displayed request inside the directed modulus lattice before a stronger request is consumed.
\item $A$ is the terminal $\RealUp$ seal row, available only after $M,S,R,D,J$ have all been displayed.
\item $H$ records componentwise $\hsame$ transport for the modulus request, stream window, regular readback, dyadic ledger, lattice comparison, seal, replay, provenance, and local name rows.
\item $C$ records displayed $\Cont$ replay from a modulus request through $M,S,R,D,J,A$.
\item $P$ is the $\Pkg$ provenance over $\BishopCauchyModulusEnvelopeUp$, $\StreamNameUp$, $\RegSeqRatUp$, $\DyadicRatCoreUp$, $\DyadicCauchyModulusLatticeUp$, $\RealUp$, $\BHist$, $\hsame$, $\Cont$, and $\NameCert$.
\item $N$ is the local $\NameCert_{\BishopCauchyModulusEnvelopeUp}$ row.
\end{enumerate}
The carrier has no coordinate for quotient representatives, selected limits, countable choice, a host modulus oracle, ambient $\RealUp$ completeness, Lean verification, or bridge checking.
\end{definition}

\begin{theorem}[Bishop Cauchy modulus envelope NameCert obligations]
\label{thm:bishop-cauchy-modulus-envelope-namecert-obligations}
For every accepted $\BishopCauchyModulusEnvelopeUp$ carrier $E=(M,S,R,D,J,A,H,C,P,N)$, the local $\NameCert_{\BishopCauchyModulusEnvelopeUp}$ surface consists exactly of carrier admission for the modulus request row, $\StreamNameUp$ finite-window exposure, $\RegSeqRatUp$ readback, $\DyadicRatCoreUp$ tolerance ledger, $\DyadicCauchyModulusLatticeUp$ comparison, $\RealUp$ seal nonescape, componentwise $\hsame$ transport, displayed $\Cont$ replay, $\Pkg$ provenance, and local naming.
\end{theorem}
\begin{proof}
Project the carrier of \autoref{def:bishop-cauchy-modulus-envelope-carrier}. The row $M$ is the only source of requested precision, and the row $S$ opens only the finite windows demanded by that request.

The readback row $R$ and dyadic ledger $D$ interpret the same displayed window. The lattice row $J$ then compares the request with finite directed-modulus data before any stronger completion-facing request is consumed.

The terminal row $A$ is downstream of $M,S,R,D,J$. The structural rows $H,C,P,N$ transport, replay, source, and name that same finite route, so the NameCert obligations are exactly the displayed rows and no hidden completion principle is available.
\end{proof}

\begin{theorem}[Bishop Cauchy modulus envelope Real-seal boundary]
\label{thm:bishop-cauchy-modulus-envelope-real-seal-boundary}
Every accepted $\BishopCauchyModulusEnvelopeUp$ packet exposes its $\RealUp$ seal only through the finite route
$$
\StreamNameUp \longrightarrow \RegSeqRatUp \longrightarrow \DyadicRatCoreUp \longrightarrow \DyadicCauchyModulusLatticeUp \longrightarrow \RealUp .
$$
Thus the envelope ties a Bishop Cauchy modulus request to Real and Completion consumers without importing quotient equality, selected limits, countable choice, a host modulus oracle, or ambient $\RealUp$ completeness.
\end{theorem}
\begin{proof}
Use \autoref{thm:bishop-cauchy-modulus-envelope-namecert-obligations}. The accepted carrier displays the modulus request before any finite stream window is opened, and the regular rational readback and dyadic comparisons remain attached to that request.

The comparison row $J$ is the only component that relates the request to a dyadic Cauchy modulus lattice. Any stronger tolerance request must therefore pass through the displayed directed lattice row rather than through an external modulus function.

The seal row $A$ is exposed only after this finite route has been replayed. Since the carrier has no coordinate for quotient equality, selected limits, countable choice, host modulus data, or ambient completeness, the Real-facing read cannot leave the displayed envelope.
\end{proof}

\falsifiablePrediction{A $\BishopCauchyModulusEnvelopeUp$ packet fails this NameCert surface if a Real or Completion consumer can read the terminal seal without first displaying the modulus request, StreamName finite windows, RegSeqRat readback, DyadicRatCore tolerance ledger, DyadicCauchyModulusLattice comparison row, transport, replay, provenance, and local naming.}

\independenceWitness{streamname, regseqrat, dyadicratcore, dyadic-cauchy-modulus-lattice, real: no listed sibling carries the whole envelope $E=(M,S,R,D,J,A,H,C,P,N)$, because each sibling supplies only finite windows, readback, tolerance arithmetic, directed modulus comparison, or terminal Real sealing while this packet binds the request-to-seal route into one local NameCert.}

\closureat{BishopCauchyModulusEnvelopeUp}{seedStr}

\begin{closurestatus}{\BishopCauchyModulusEnvelopeUp}
  \constructivestory{A $\BishopCauchyModulusEnvelopeUp$ packet is generated from a modulus request row, StreamName finite windows, RegSeqRat readback, a DyadicRatCore tolerance ledger, a DyadicCauchyModulusLattice comparison row, a RealUp terminal seal, componentwise hsame transport, displayed Cont replay, Pkg provenance, and a local NameCert row.}
  \theoryclosure{\seedClosure}
  \scopeclosed{\autoref{def:bishop-cauchy-modulus-envelope-carrier}, \autoref{thm:bishop-cauchy-modulus-envelope-namecert-obligations}, and \autoref{thm:bishop-cauchy-modulus-envelope-real-seal-boundary} seed the finite Bishop Cauchy modulus envelope over StreamName, RegSeqRat, DyadicRatCore, DyadicCauchyModulusLattice, and Real rows.}
  \formalstatus{\unformalizedV}
  \bridgestatus{none}
  \notclaimed{This chapter does not close quotient real equality, selected limits, countable choice, host modulus functions, ambient Real completeness, Lean verification, or bridge checking.}
  \upgradepath{Obligation closure requires checked carrier admission, modulus request exposure, finite-window readback, dyadic tolerance routing, lattice comparison, Real seal nonescape, transport, replay, provenance, and local naming for BEDC.Derived.BishopCauchyModulusEnvelopeUp.}
\end{closurestatus}
